%! Rational Functions and End Behavior — Unit 3, Lesson 3.1 — Group Activity (Answer Key)
\documentclass[10pt]{article}
\usepackage{precalculus-article}
\usepackage{precalculus-key}   % pulls in -boxes; do NOT also load -boxes
\newcommand{\TallMath}[1]{$\displaystyle #1\rule[-1.4em]{0pt}{3.2em}$}

\begin{document}

\pageheader{Unit 3, Lesson 3.1}{Group Activity --- Answer Key}
\namepartnerperiod

\vspace{0.10in}

\begin{scenariobox}[Analyzing Concert Venue Costs]{sky}
\small
A concert venue is modeling the average cost per person for different events.
Each rational function below represents a different cost or demand model.
Analyze the end behavior of each function as audience size $x$ grows without bound.

\medskip
\renewcommand{\arraystretch}{1.5}
\begin{tabular}{@{}llll@{}}
$r_1(x) = \dfrac{2x + 5}{x - 1}$ &\quad&
$r_2(x) = \dfrac{3x^2}{x^2 + 4}$ \\[8pt]
$r_3(x) = \dfrac{x + 1}{x^2 - 9}$ &\quad&
$r_4(x) = \dfrac{x^3 - 2x}{x - 3}$ \\
\end{tabular}
\end{scenariobox}

\vspace{0.08in}

\begin{tcolorbox}[colback=white, colframe=black!40,
    title=\textbf{Tier R --- Remediate},
    fonttitle=\bfseries, arc=2mm, left=3mm, right=3mm, top=2mm, bottom=2mm]

The leading terms are already extracted for you.
Compute the quotient of leading terms and state the asymptote type.

\medskip
\renewcommand{\arraystretch}{2.0}
\begin{tabular}{@{}p{1.2cm} p{2.2cm} p{2.2cm} p{2.8cm} p{3.0cm}@{}}
\textbf{Fcn} & \textbf{Lead.\ num} & \textbf{Lead.\ den} & \textbf{Quotient} & \textbf{Asymptote?} \\
\hline
$r_1$ & $2x$   & $x$   & \ans{$2x/x = 2$}    & \ans{HA: $y=2$} \\
$r_2$ & $3x^2$ & $x^2$ & \ans{$3x^2/x^2 = 3$} & \ans{HA: $y=3$} \\
$r_3$ & $x$    & $x^2$ & \ans{$x/x^2 = 1/x \to 0$} & \ans{HA: $y=0$} \\
$r_4$ & $x^3$  & $x$   & \ans{$x^3/x = x^2$}  & \ans{Polynomial (no HA)} \\
\end{tabular}
\end{tcolorbox}

\vspace{0.08in}

\begin{tcolorbox}[colback=white, colframe=black!40,
    title=\textbf{Tier A --- Apply},
    fonttitle=\bfseries, arc=2mm, left=3mm, right=3mm, top=2mm, bottom=2mm]

For each function, identify the degrees and leading coefficients, compute the quotient of leading
terms, write both limit statements, and state the asymptote type.

\medskip
\renewcommand{\arraystretch}{2.2}
\begin{tabular}{@{}p{1.2cm} p{1.2cm} p{1.2cm} p{2.5cm} p{2.8cm} p{2.8cm} p{1.8cm}@{}}
\textbf{Fcn} & \textbf{Deg n} & \textbf{Deg d} & \textbf{Quotient} &
$\displaystyle\lim_{x\to\infty}$ & $\displaystyle\lim_{x\to -\infty}$ & \textbf{Type} \\
\hline
$r_1$ & \ans{1} & \ans{1} & \ans{$2$}     & \ans{$2$}      & \ans{$2$}      & \ans{HA} \\
$r_2$ & \ans{2} & \ans{2} & \ans{$3$}     & \ans{$3$}      & \ans{$3$}      & \ans{HA} \\
$r_3$ & \ans{1} & \ans{2} & \ans{$0$}     & \ans{$0$}      & \ans{$0$}      & \ans{HA} \\
$r_4$ & \ans{3} & \ans{1} & \ans{$x^2$}   & \ans{$+\infty$} & \ans{$+\infty$} & \ans{Poly} \\
\end{tabular}
\end{tcolorbox}

\vspace{0.08in}

\begin{tcolorbox}[colback=white, colframe=black!40,
    title=\textbf{Tier E --- Extend},
    fonttitle=\bfseries, arc=2mm, left=3mm, right=3mm, top=2mm, bottom=2mm]

\begin{enumerate}[itemsep=10pt]
  \item \textbf{Construct a rational function} with a \textbf{horizontal asymptote at $y = -3$}
        where the numerator has degree 2 and the denominator has degree 2.

        Write your function: $r(x) = $ \ans{$\dfrac{-3x^2 + 1}{x^2 + 5}$} \textit{(many answers possible)}

        Verify: What are the leading terms?  Quotient $=$ \ans{$-3x^2/x^2 = -3$}. Does it equal $-3$? \ans{Yes}

  \item \textbf{Construct a rational function} with a \textbf{slant asymptote at $y = 2x + 1$}.

        Write your function: $r(x) = $ \ans{$\dfrac{2x^2 + x + 3}{x}$} \textit{(many answers possible)}

        \ansline{The numerator has degree 2 and the denominator has degree 1, so deg num $=$ deg denom $+ 1$.
        This means the quotient of leading terms is a linear expression ($2x$), giving a slant asymptote
        rather than a constant, so no horizontal asymptote exists.}

  \item For your function in part (1), write both limit statements:

        $\displaystyle\lim_{x\to\infty} r(x) = $ \ans{$-3$}
        \qquad
        $\displaystyle\lim_{x\to -\infty} r(x) = $ \ans{$-3$}
\end{enumerate}
\end{tcolorbox}

\end{document}
